\section{附录}
\label{sec:appendix}

\begin{table*}[!ht]
	\centering
	\begin{tabular}{lllllll}
		\toprule
		\textbf{实例类型} & \textbf{加速器} & \makecell{\textbf{本地 SSD}\\\textbf{带宽/GPU}} & \makecell{\textbf{远程 SSD}\\\textbf{带宽/GPU}} & \makecell{\textbf{网络}\\\textbf{带宽/GPU}} & \textbf{含 NVLink} & \textbf{价格} \\ \midrule
		a2-ultragpu-8g \cite{googleinstance}      & 8 x A100(80\,GB) & 2.58\,Gbps & 0.29\,Gbps & 12.5\,Gbps & \ding{51} & 40.44 USD/h \\
		p4d.24xlarge\cite{awsinstance}            & 8 x A100(40\,GB) & 2.31\,Gbps & -          & 100\,Gbps  & \ding{51} & 45.039 USD/h \\
		ml.hpcpni2.28xlarge\cite{volcanoinstance} & 8 x A100(80\,GB) & 4\,Gbps    & -          & 100\,Gbps  & \ding{55} & 48.23 USD/h \\
		p4de.24xlarge\cite{awsinstance}           & 8 x A100(80\,GB) & 2.31\,Gbps & -          & 100\,Gbps  & \ding{51} & 56.328 USD/h \\
		a3-highgpu-8g\cite{googleinstance}        & 8 x H100         & 6.09\,Gbps & 0.97\,Gbps & 100\,Gbps & \ding{51} & 88.25 USD/h \\
		a3-megagpu-8g\cite{googleinstance}        & 8 x H100         & 6.09\,Gbps & 0.97\,Gbps & 200\,Gbps & \ding{51} & 不可用 \\
		p5.48xlarge \cite{awsinstance}            & 8 x H100         & 9.8\,Gbps  & -          & 400\,Gbps  & \ding{51} & 不可用 \\ \bottomrule
	\end{tabular} \\[10pt]

    \begin{minipage}{1\linewidth}
        \caption{\small{\textit{
		来自 GPU 云厂商的 {\maas} 硬件配置概览。
        }}}
    \label{tab:vendor-hardware}
    \end{minipage}
\end{table*}

\subsection{值得说明的实现细节}
\label{sec:appendix-implementation}

\nospacestitle{网络库。 }
我们实现了一个通信库，
统一抽象 NVLink 和 RDMA，
从而整体性地完成参数传输，
这一点与 DeepEP~\cite{deepep} 类似。
在实现过程中我们发现，
如果直接使用现成的组通信器（例如 NCCL~\cite{nccl}），
跨机器建立通信组的开销较大（例如 100\,ms），
这会显著限制基于网络扩缩容的效果。
幸运的是，我们的计划只需要节点对之间的 P2P 通信。
因此，我们预先构建了一个连接池，
支持各节点之间的全互连连接。
尽管计算网络（RDMA）在可扩展性上存在潜在问题~\cite{fasst}，
但这主要出现在小负载传输场景中，
并且可以通过 DCT~\cite{krcore} 等先进 RDMA 传输机制来缓解。

\stitle{带 CUDA 上下文池的原生服务引擎。 }
在执行之前，GPU 上必须先创建一个带已加载内核（cuModule）的 CUDA 上下文。
创建这样一个 CUDA 上下文大约需要 500\,ms，
在服务实例自动扩缩容中不可忽略。
为降低该开销，{\sys} 维护了一个小规模 CUDA 上下文池，
其中预先加载好常用内核，
然后在已有 CUDA 上下文中把参数传入 GPU，
这一做法与已有工作~\cite{DBLP:journals/corr/abs-2405-12079} 类似。
此外，{\sys} 使用 C++ 与原生 CUDA API 构建，
从而消除了初始化 PyTorch（例如 \texttt{dlopen}）带来的额外开销。

\stitle{容错。 }
发生机器故障时，
我们会利用自身的扩缩容机制自动扩出新的实例。
其中一个问题是，故障机器上的缓存参数会一并丢失，
因此需要把这些参数重新分发到其他机器上，
以维持全局参数池的不变式。
对于调度器、监控器等其他组件的故障，
我们沿用现有工作中的恢复流程~\cite{DBLP:conf/usenix/Romero0YK21,DBLP:journals/corr/abs-2401-14351}。

\subsection{{\maas} 的硬件配置}
\label{sec:appendix-hardware}

\noindent
表~\ref{tab:vendor-hardware} 列出了典型 {\maas} 系统所依赖的硬件配置。
